%% 仿射 Grassmann 簇
\section[仿射 Grassmann 簇]{仿射 Grassmann 簇（4 月 29 日）}

\srcpage{71}

今天我们要讨论几何 Satake 对应，虽然有点简短。TODO：我今天比平时更困惑，所以如果有什么说不通，很可能是我写错了。对此抱歉。

我们一直沿着一个总主题：若 $G$ 是作用在空间 $X$ 上的群，$K$ 是 $G$ 的子群，则商 $X/K$ 带有群胚作用。这导致 Hecke 代数作用在诸如 $X/K$ 上的函数之类的东西上。

上次我们引入了一个对应
\[
  \begin{tikzcd}
    & \Hck_x \arrow[dl] \arrow[dr] & \\
    \Bun_G(C) & & \Bun_G(C),
  \end{tikzcd}
\]

\srcpage{72}

其中 $C$ 是曲线、$x\in C$。左映射是一个纤维丛，其纤维为
\[
  \Gr_G := G(K_x)/G(\cO_x).
\]
这个空间称为 $G$ 的仿射 Grassmann 簇（affine Grassmannian）。空间 $\Hck_x$ 是 $C$ 上主丛对的空间，附带 $C\setminus x$ 上的一个等同数据。作为 $\Bun_G(C)$ 上的纤维丛，$\Hck_x$ 是某个 $LG$-丛的配丛；这里 $LG=G(K_x)$，$LG_+=G(\cO_x)$；它们分别扮演 $G$ 与 $K$ 的角色。这个故事里重要的是 Beauville–Laszlo 定理（TODO：拼写对吗？），它告诉我们如何看待圆盘上的 $G$-丛。

你可以把 $\Gr_G$ 想成分类下列东西：$C$ 上带 $C\setminus x$ 上平凡化的主 $G$-丛；等价地，圆盘上带原点之外平凡化的主 $G$-丛。给定 $P\in\Bun_G(C)$，我们可以在一个圆盘上把它平凡化，得到 $P$ 及其平凡化，即 $\Bun_G(C,x)$ 中的一个元素。然后修正圆盘上的平凡丛，得到 $(P,P')\in\Bun_G(C,x)\times\Gr_G$。现在模去 $G(\cO_x)$ 的对角作用——它改变 $P$ 与 $P'$ 的平凡化——这个商就是仿射 Grassmann 簇，并把 $\Gr_G$ 实现为一个配丛。

直接指定 $\Hck_x$ 中的点很难：局部看它像 $\Gr_G\times\Bun_G(C)$，但整体上是一个扭乘积。你可以用 TODO 以非不变的方式指定点。令 $\mathbb{O}\subset\Gr_G$ 是一个陪集，等价地是 $LG_-\backslash LG/LG_+$ 的一个元素。于是 $\mathbb{O}$ 定义了 $\Hck_x$ 的一个子集，而事实上对应保持这个子集。也就是说，我们有图
\[
  \begin{tikzcd}
    & \mathbb{O} \arrow[dl] \arrow[dr] & \\
    \Bun_G(C) & & \Bun_G(C).
  \end{tikzcd}
\]
所以从某种意义上说，这些轨道给了我们许多可能更小的对应操作。这很好，因为 $\Hck_x$ 里的群运算并不多，而这些轨道给了我们更多选择。

\begin{example}
令 $G=\mathrm{GL}_n$。具体地，$\Gr_{\mathrm{GL}_n}$ 是圆盘上平凡秩 $n$ 向量丛的修正：我们取 $V\in\Vect(D)$，连同 $V$ 在去心圆盘 $D^\times$ 上的一个平凡化。于是 $\Gr_{\mathrm{GL}_n}$ 含有一个子空间，其中 $V$ 在整个 $D$ 上可平凡化；这给出 $V(D^\times)\cong K^n$，即 Laurent 级数（TODO：这完全说不通，但我想这就是我听到的？抱歉。）。TODO：关于格的某些话。

$V(D)$ 中的极点是有界的：$t^{-N}\cO^n\subset V(D)\subset t^N\cO^n\subset K^n$。这说明了如何把 $\Gr_{\mathrm{GL}_n}$ 实现为（TODO：某物）中射影簇之并。因此 $\Gr_{\mathrm{GL}_1}$ 是 ind-紧合 ind-概形。对更一般的约化群 $G$，我们可以把 $G\hookrightarrow\mathrm{GL}_n$ 嵌入某个 $n$，于是 $\Gr_{\mathrm{GL}_n}$ 的这些好性质推广到 $\Gr_G$。\yizhu{原文此处作 $\mathrm{Gr}_{\mathrm{GL}_1}$，据上下文疑为 $\mathrm{Gr}_{\mathrm{GL}_n}$ 之误。译文照原文排印。}{A（原文 p.72）}
\end{example}

任何在 $G(\cO)$ 下不变的几何，在 $\Bun_G(C)$ 上都有意义。例如，$\Gr_G$ 上 $G(\cO)$-等变层下降为 $\Hck_x$ 上的层。这包括轨道上的常层。又因为 $\Hck_x$ 是 $\Bun_G(C)$ 上的群胚，我们得到一个从 $\Hck_x$ 上的层到 $\Bun_G(C)$ 上层的自同态的映射。特别地，$\Gr_G$ 上 $G(\cO)$-等变层的范畴，等价于 $G(\cO)\backslash G(K)/G(\cO)$ 上层的范畴；上次我们称之为球面 Hecke 范畴。特别地，球面 Hecke 范畴幺半地作用在 $\Bun_G(C)$ 上层的范畴上。

上次我们论证了：球面 Hecke 范畴上的幺半结构实际上是对称幺半的。办法是场论论证：有三个方向可以移动算子，足以容纳一个对称幺半结构。

由此可以提取出一种直接证明交换性的办法：设法看清「把这些局部算子绕来绕去」是什么意思。这正是 Beilinson–Drinfeld 用因子化代数理论所做的。这用到了 $\Gr_G$ 上一些额外的代数结构，涉及变动点 $x\in C$。

另一个视角来自同伦论：当 $X$ 是空间时，$\pi_2(X)$ 是阿贝尔群，原因大体类似。所以让我们深入 $\Gr_G$ 的同伦论。在阿贝尔情形，我们得到了同伦等价于 $\Z$ 的东西。一般地，$\Gr_G$ 可以用 $\A^1$ 上带有 $\A^1\setminus 0$ 上平凡化的 $G$-丛来建模。由于本段只关心同伦，我们可以把 $\A^1$ 换成 $\R^2$，把 $\A^1\setminus 0$ 换成 $\R^2\setminus D_\varepsilon(0)$：把原点的形式邻域换成真正的邻域。接着把 $G$ 换成它的极大紧子群 $G_c$：我们看的是 $\R^2$ 上、在小圆盘外有平凡化的主 $G_c$-丛。主 $G_c$-丛意味着到 $BG_c$ 的映射，所以在同伦意义下，
\[
  \Gr_G \simeq \mathrm{Map}((\R^2,\R^2\setminus D_\varepsilon(0)),(BG_c,*))
  = \Omega^2 BG_c = \Omega BG_c.
\]

\srcpage{73}

包含 $G_c\hookrightarrow G$ 总是诱导同伦等价 $BG_c\xrightarrow{\sim}BG$。无论如何，仿射 Grassmann 簇的这个描述对拓扑信息有用：例如 $\pi_0(\Gr_G)=\pi_1(G)=\pi_2(BG)$。例如 $\Gr_{\mathrm{PGL}_n}$ 连通，而 $\Gr_{\mathrm{SL}_n}$ 不连通。

一般地，在同伦论中，把一个空间实现为别的空间的 $\Omega^k$，是提供部分交换乘法的一种方式：环路可以复合，但有 $k$ 个可以复合的方向。所以既然 $\Gr_G$ 在同伦意义下是某个东西的 $\Omega^2$，$\Gr_G$ 就（同伦意义下）有一个「弱交换」的群结构。

同伦论者用小球 $k$-盘操作子（little $k$-discs operad），或 $E_k$ 操作子，来形式化 $\Omega^kX$ 上「弱交换乘法」的想法。这是一个操作子，其上在 $\mathbf{Top}$ 中的代数恰好是 $k$ 重环路空间。也就是说，空间上的 $E_k$-结构给你这样的能力：给定更大的 $D^k$ 中 $m$ 份 $D^k$，或 $\mathrm{Conf}_m(D^k)$ 中的一个点，我们就得到 $X$ 上的一个乘法。\yizhu{$E_k$ 操作子（operad）是现代同伦论的标准工具：它编码「有 $k$ 个可交换方向」的乘法，$E_\infty$ 才给出严格交换。}{I（术语表 operad＝操作子）}

事实上，$\pi_k(X)$ 带有 $E_k$-结构，这里我们把从圆盘出发的映射想成把边界送到基点。于是按算子乘积展开方式的复合，与 $\pi_k(X)$ 中的复合相同。对 $k\ge 2$，集合上的 $E_k$-结构就是阿贝尔群结构；所以 $k\ge 2$ 时 $\pi_k$ 是阿贝尔的。

如上所见，$\Gr_G$ 有 $E_2$-结构。所以无论以何种方式线性化它（函数、层……），我们都会得到 $E_2$-对象。例如，仿射 Grassmann 簇上的层是一个 $E_2$-范畴，即「弱交换」幺半（实际上是编织幺半）。这正是在 3 维中线算子范畴上得到的结构。

$\Gr_G$ 上的 $E_2$-结构可以想成由在 $C$ 上沿两个不同方向移动 $x\in C$ 给出，而 $C$ 是二维的。但 $G(\cO_x)\backslash\Gr_G$ 中还多一个交换性的维度（你也许已经预期到：我们在 4 维理论中有线算子，而不是 3 维，所以预期多一点交换性）。除了来自 $\Gr_G$ 的 $E_2$-结构之外还有复合，于是得到 $E_3$-结构，正如 4 维中线算子所应有。

\begin{remark}
仿射 Grassmann 簇是 ind-概形（ind-scheme），因此我们必须担心它是否既约。当 $G$ 半单时，$\Gr_G$ 既约；但即便对 $\mathrm{GL}_1$，仿射 Grassmann 簇其实是 $\Z$ 乘以一个无穷维形式群！这对我们正在做的事（拓扑：取可构造层或 D-模之类）不重要。但它对切空间确实重要：$T\Gr_G=K/O$；例如对 $\mathrm{GL}_1$，我们得到 $\C((t))/\C[[t]]$，这是个无穷维 Lie 代数！取指数，我们得到一个无穷维形式群。

所以这是看出非既约性的一种方式。无论如何，我们主要会坚持半单群。
\end{remark}

Pressley–Segal（TODO：引用）是关于仿射 Grassmann 簇的另一个好参考，视角不同。

$\Gr_G$ 上有 Białynicki-Birula 给出的 Morse 理论视角，给出以单参数子群 $\G_m\to T$ 为指标的一个胞腔分解，其中 $T$ 是 $G$ 的极大环面；这些单参数子群称为余权（coweight）。余权构成一个格，记作 $\Lambda$。这个胞腔分解是 Schubert 分解的无穷维版本。$\Gr_G$ 上有由旋转（去心）圆盘给出的 $\G_m$-作用。

我们也可以选一个一般的余权 $\rho^\vee$（意思是其像的中心化子就是 $T$ 本身）；由于 $G\subset LG_+$ 作用在 $\Gr_G$ 上，$\rho^\vee$ 定义了 $\Gr_G$ 上另一个 $\G_m$-作用。

让我们把这些作用合起来，寻找不动点。

它们由余特征 $\phi:\G_m\to G$ 给出，这些映射被送到（TODO）。例如对 $G=\mathrm{GL}_n$，我们要找对角矩阵，对角元为 $(t^{i_1},\dots,t^{i_n})$，它是 $\mathrm{GL}_n(\C((t)))$ 的一个元素。这是一个自然的出发点：一般的 Morse 理论论证说，孤立临界点给出一个胞腔分解，每个胞腔中心有一个不动点（胞腔就是稳定流形）。

余特征格 $\Lambda$（TODO：这不是余权的记号吗？）同构于 $T(K)/T(O)\cong\Gr_T$。这很好，因为我们理解环面的仿射 Grassmann 簇：$\Lambda\cong\Gr_T$，而它单射进 $\Gr_G$。我们将用这一点，借助阿贝尔数据来理解 $\Gr_G$ 上层的范畴。下一步是找出 Cartan 分解；这与按 $G(O)$-轨道的分解并不完全相同，但每个 $G(O)$-轨道都是胞腔之并。事实上 $G(O)\backslash\Gr_G\cong\Lambda/W$，其中 $W$ 是 Weyl 群。这个事实的证明是线性代数。所以我们关心的胞腔并不恰好与余特征一一对应，而是与模去 Weyl 群作用的余特征一一对应。

下次我们将把这个视角与交换乘法结合起来，得到几何 Satake 定理。

\srcpage{74}
